% =====================================================================
\chapter{Precisión numérica y precisión mixta}
\label{tema:4}
% =====================================================================

Los núcleos tensoriales son mucho más rápidos en 16 u 8 bits que en 32, y cada bit que se ahorra reduce también la memoria y el tráfico. El precio es un mayor error de redondeo y un rango más estrecho. Este tema formaliza los formatos de coma flotante, acota el error de redondeo y de la suma acumulada, y justifica con esas cotas cada elemento de la receta de entrenamiento en precisión mixta. Se siguen los capítulos~2 a~4 de Higham~\cite{higham}, el artículo de Micikevicius \emph{et al.}~\cite[§3]{micikevicius}, la propuesta de formatos fp8~\cite{fp8}, las clases de cuantización de MIT 6.5940~\cite{mit65940} y el §2.1 de la tarea~2 de CS336~\cite{cs336a2}. El módulo~11 introdujo estos formatos de forma descriptiva; aquí se demuestran sus propiedades.

\section{Formatos de coma flotante}

\begin{definicion}[Formato binario de coma flotante]\label{def:formato}
Un formato binario con $e$ bits de exponente y $p$ bits de mantisa representa los números \emph{normales}
\[
x=(-1)^{\sigma}\,(1{,}f_1f_2\dots f_p)_2\cdot 2^{E},\qquad E_{\min}\le E\le E_{\max},
\]
con $\sigma\in\{0,1\}$ y $f_i\in\{0,1\}$. En los formatos IEEE~754 el exponente se guarda con un sesgo $b=2^{e-1}-1$, y $E_{\min}=1-b$, $E_{\max}=b$; los códigos de exponente extremos se reservan para el cero, los números subnormales, los infinitos y los NaN.
\end{definicion}

\begin{proposicion}[Rango y separación]\label{prop:rango}
En el formato de la \cref{def:formato}:
\begin{enumerate}
  \item el mayor número representable es $(2-2^{-p})\,2^{E_{\max}}$ y el menor normal positivo, $2^{E_{\min}}$;
  \item en cada intervalo $[2^{E},2^{E+1})$ con $E_{\min}\le E\le E_{\max}$ hay $2^{p}$ números representables equiespaciados, separados por $\ulp=2^{E-p}$;
  \item la separación entre $1$ y el siguiente representable es el \emph{épsilon de la máquina}, $\varepsilon=2^{-p}$.
\end{enumerate}
\end{proposicion}

\begin{table}[h]
\centering\small
\begin{tabular}{@{}lccccc@{}}
\toprule
Formato & bits (s-e-m) & máximo & menor normal & $\varepsilon=2^{-p}$ & uso típico \\
\midrule
fp32 & 1-8-23 & $3{,}40\cdot 10^{38}$ & $1{,}18\cdot 10^{-38}$ & $1{,}19\cdot 10^{-7}$ & copia maestra, acumulación \\
fp16 & 1-5-10 & $65\,504$ & $6{,}10\cdot 10^{-5}$ & $9{,}77\cdot 10^{-4}$ & productos \\
bf16 & 1-8-7 & $3{,}39\cdot 10^{38}$ & $1{,}18\cdot 10^{-38}$ & $7{,}81\cdot 10^{-3}$ & productos, activaciones \\
fp8 E5M2 & 1-5-2 & $57\,344$ & $6{,}10\cdot 10^{-5}$ & $0{,}25$ & gradientes \\
fp8 E4M3 & 1-4-3 & $448$ & $1{,}56\cdot 10^{-2}$ & $0{,}125$ & pesos, activaciones \\
\bottomrule
\end{tabular}
\caption{Formatos de coma flotante usados en el entrenamiento. El formato E4M3 no sigue del todo IEEE~754: no tiene infinitos y reutiliza el código de exponente máximo, lo que eleva su máximo de 240 a 448~\cite{fp8}.}
\label{tab:formatos}
\end{table}

El rango depende de los bits de exponente y la precisión, de los de mantisa. bf16 conserva los 8 bits de exponente de fp32 y, por tanto, su rango; fp16 tiene 3 bits más de mantisa, pero su rango es mucho menor. Por debajo del menor normal, los números subnormales extienden el rango con pérdida de precisión: en fp16 el menor positivo es $2^{-24}\approx 5{,}96\cdot 10^{-8}$.

\section{Error de redondeo}

\begin{definicion}[Redondeo y unidad de redondeo]
Sea $\fl(x)$ el número representable más cercano a $x\in\R$ (redondeo al más cercano, con empates al par). La \emph{unidad de redondeo} del formato es $u=\varepsilon/2=2^{-p-1}$.
\end{definicion}

\begin{teorema}[Error relativo de redondeo]\label{teo:redondeo}
Si $x$ está en el rango normal del formato, $\fl(x)=x(1+\delta)$ con $\abs{\delta}\le u$.
\end{teorema}

La norma IEEE~754 exige que las operaciones básicas se redondeen correctamente: $\fl(a\circ b)=(a\circ b)(1+\delta)$, $\abs{\delta}\le u$, para $\circ\in\{+,-,\times,/\}$~\cite[cap.~2]{higham}. Es el \emph{modelo estándar} de la aritmética en coma flotante.

\begin{lema}[Estancamiento]\label{lem:estancamiento}
Sea $s>0$ representable con $2^{E}\le s<2^{E+1}$ y sea $0\le a<2^{E-p-1}=\ulp(s)/2$. Entonces $\fl(s+a)=s$.
\end{lema}

\begin{corolario}[Actualizaciones perdidas]\label{cor:actualizacion}
En bf16, si un peso vale $w=1$, toda actualización $0\le\Delta w<2^{-8}\approx 0{,}0039$ se pierde: $\fl(1+\Delta w)=1$.
\end{corolario}

\begin{lema}[Producto de factores de redondeo]\label{lem:gamma}
Sean $\delta_1,\dots,\delta_k$ con $\abs{\delta_i}\le u$ y $ku<1$. Entonces
\[
\Bigl\lvert\prod_{i=1}^{k}(1+\delta_i)-1\Bigr\rvert\le\gamma_k,\qquad\gamma_k=\frac{ku}{1-ku}.
\]
\end{lema}

\begin{teorema}[Error de la suma recursiva]\label{teo:suma}
Sea $\hat s_n$ el resultado de sumar $x_1,\dots,x_n$ en coma flotante en el orden $\hat s_1=x_1$, $\hat s_k=\fl(\hat s_{k-1}+x_k)$. Si $(n-1)u<1$ y no hay desbordamiento,
\[
\abs{\hat s_n-s_n}\le\gamma_{n-1}\sum_{i=1}^{n}\abs{x_i},\qquad s_n=\sum_{i=1}^{n}x_i.
\]
\end{teorema}

La cota crece con $n$. En fp16, $u=2^{-11}\approx 4{,}9\cdot 10^{-4}$ y con $n=1000$ sumandos $\gamma_{999}\approx 0{,}95$: la cota no garantiza ni una cifra. En fp32, $u=2^{-24}\approx 6\cdot 10^{-8}$ y $\gamma_{999}\approx 6\cdot 10^{-5}$. Esta es la razón por la que los núcleos tensoriales acumulan los productos escalares en fp32 aunque los factores estén en 16 bits (\cref{eje:4-acumular}).

\section{La receta de la precisión mixta}

Micikevicius \emph{et al.}~\cite[§3]{micikevicius} proponen tres técnicas, que siguen siendo la base de la práctica actual.

\begin{enumerate}[label=(\arabic*)]
  \item \emph{Copia maestra de los pesos en fp32} (§3.1). Los productos usan una copia en 16 bits, pero la actualización $w\gets w-\eta g$ se aplica a una copia en fp32. Por el \cref{cor:actualizacion}, una actualización pequeña frente al peso se perdería en 16 bits; los autores observan que esto ocurre cuando el peso es unas $2^{11}$ veces mayor que la actualización.
  \item \emph{Escalado de la pérdida} (§3.2). En fp16, los valores menores que $2^{-24}$ se anulan, y parte de los gradientes de activación caen por debajo. Se multiplica la pérdida por una constante $c$ antes de la retropropagación y se dividen los gradientes por $c$ antes de actualizar.
  \item \emph{Acumulación en fp32} (§3.3). Los productos escalares acumulan en fp32, y las reducciones grandes (softmax, normalizaciones, pérdida) se calculan en fp32.
\end{enumerate}

\begin{proposicion}[El escalado por potencias de 2 es exacto]\label{prop:escalado}
Sea $c=2^{k}$ con $k\in\mathbb{Z}$, y sea $x$ tal que $x$ y $cx$ están en el rango normal. Entonces $\fl(cx)=c\,\fl(x)$. En particular, multiplicar y después dividir por $c$ un número representable no introduce error.
\end{proposicion}

Como la retropropagación es lineal en la semilla $\partial\mathcal{L}/\partial\mathcal{L}=1$, multiplicar la pérdida por $c$ multiplica por $c$ todos los gradientes, y por la \cref{prop:escalado} la división final por $c=2^{k}$ los restituye sin error adicional, salvo en los valores que se hayan desbordado o anulado. En la práctica se usa un escalado \emph{dinámico}: se duplica $c$ periódicamente y se reduce a la mitad, descartando el paso, cuando aparece un infinito. Con bf16, cuyo rango es el de fp32, el escalado no es necesario, y esa es la razón principal de que bf16 sea hoy el formato habitual de entrenamiento.

\begin{observacion}[Memoria de la receta]
Con Adam y precisión mixta, cada parámetro ocupa 2~bytes (pesos en 16 bits) $+$ 2 (gradiente en 16 bits) $+$ 12 (copia maestra, primer y segundo momento en fp32) $=16$~bytes~\cite[§3.1]{zero}. El \cref{tema:5} estudia cómo repartir esos 16~bytes entre dispositivos.
\end{observacion}

\section{Cuantización}

El entrenamiento en fp8 y la inferencia en int8 o int4 llevan la idea al extremo. Con tan pocos bits de mantisa, cada tensor necesita su propio factor de escala.

\begin{definicion}[Cuantización lineal simétrica]\label{def:cuantizacion}
Sea $x\in\R^{L}$ y $b\ge 2$ un número de bits. La \emph{cuantización lineal simétrica} de $x$ con $b$ bits usa el paso $\Delta=\max_i\abs{x_i}/(2^{b-1}-1)$ y los enteros $q_i=\operatorname{round}(x_i/\Delta)\in\{-(2^{b-1}-1),\dots,2^{b-1}-1\}$. La \emph{descuantización} es $\hat x_i=\Delta\,q_i$.
\end{definicion}

\begin{proposicion}\label{prop:cuantizacion}
Con la \cref{def:cuantizacion}, $\abs{\hat x_i-x_i}\le\Delta/2$ para todo $i$.
\end{proposicion}

El error depende del valor máximo del tensor: un solo valor atípico grande aumenta $\Delta$ y degrada la representación de todos los demás. Por eso se cuantiza por canales o por bloques, y se tratan aparte los valores atípicos~\cite{mit65940}. El módulo~19 retoma la cuantización para la inferencia.

